\Needspace{12\baselineskip}
\section{Bott--Chern classes and closed currents}\label{sec:bc-closed-currents}

Bott--Chern cohomology separates the smooth background of a closed current from its global potential. Subtracting the trace current from a closed background then gives a residual class on each model. Pushforward compatibility organizes these classes into a b-divisor class, as in \cite[section \emph{The notion of b-divisors}]{SGPT}.

\index{Bott--Chern class!definition}For a complex manifold $Z$, let $\Lambda^{1,1}_{\mathbb R}$ denote the bundle of real $(1,1)$-forms and define
\[
 H^{1,1}_{\mathrm{BC}}(Z,\mathbb R)
 =\frac{\{\beta\in C^\infty(Z,\Lambda^{1,1}_{\mathbb R}):\mathrm{d}\beta=0\}}
        {\{\ddc f:f\in C^\infty(Z,\mathbb R)\}}.
\]
\Needspace{8\baselineskip}
\begin{lemma}\label{lem:bc-one-one-current-description}
Replacing smooth forms and smooth functions in this quotient by real currents and real distributions gives the same space. In particular every closed positive $(1,1)$-current $T$ on $Z$ can be written
\[
 T=\beta+\ddc w
\]
with $\beta$ smooth and closed and $w$ a global quasi-plurisubharmonic function.
\end{lemma}
\begin{proof}
The decomposition in \cref{lem:closed-current-global-potential} proves surjectivity of the smooth-to-current map and gives a quasi-plurisubharmonic potential for positive currents. Its smooth-difference assertion also proves injectivity: a distribution whose $\ddc$ is smooth is smooth, so two smooth representatives identified in the current quotient are already identified in the smooth quotient.
\end{proof}

\index{Pseudoeffective class}For global pairings, take $Z$ connected and compact. The class of a closed current $T$ is written $\{T\}_{\mathrm{BC}}$. It is called pseudoeffective when it has a positive closed representative. Pullback of classes uses smooth representatives. Pushforward under a proper modification uses currents; it is well-defined because pushforward commutes with $\ddc$. The lemma identifies its result with a smooth Bott--Chern class. For a modification $g$, we have $g_*g^*\beta=\beta$ as currents for every smooth form $\beta$ of the indicated degree, by degree-one change of variables.

\Needspace{9\baselineskip}
A Gauduchon metric $\gamma$ defines a linear functional on this cohomology:
\[
 \ell_\gamma(\{T\}_{\mathrm{BC}})=\int_ZT\wedge\gamma^{\dim Z-1}.
\]
Current Stokes makes it independent of the representative. For a positive representative it is nonnegative, and equality forces that representative to vanish, by positivity and the strict positivity of the metric. In particular the pseudoeffective cone contains no nonzero line: if $T_1,T_2\geq0$ represent opposite classes, their sum has zero Gauduchon pairing, hence $T_1=T_2=0$.

\Needspace{12\baselineskip}
Use the Hopf manifold $X_n$, $n\geq2$, and its Gauduchon metric $\omega$ from \cref{prop:hopf-bergman-exact}, with covering coordinates $z\in\mathbb C^n\setminus\{0\}$. The form $\alpha=\ddc\log|z|^2$ is exact as a de Rham form. Indeed $\dc\log|z|^2$ descends to the quotient because the logarithm changes by a constant under the contraction. Nevertheless
\[
 \ell_\omega(\{\alpha\}_{\mathrm{BC}})
 =\int_{X_n}\alpha\wedge\omega^{n-1}=2(n-1)>0.
\]
Thus its Bott--Chern class is nonzero. Gauduchon integration detects closed-curvature information that the de Rham class loses in this example.

Return now to the smooth compact modifications of the original manifold $Y$.

\Needspace{11\baselineskip}
\begin{definition}\label{def:bott-chern-bdivisor-class}\index{Bott--Chern class!b-divisor}
A \emph{Bott--Chern b-divisor class} over $Y$ is a family
\[
 \mathbb A=(\mathbb A_Z)_{\pi:Z\to Y},\qquad
 \mathbb A_Z\in H^{1,1}_{\mathrm{BC}}(Z,\mathbb R),
\]
over smooth compact modifications, compatible with pushforward. It is pseudoeffective if every $\mathbb A_Z$ is pseudoeffective. It is \emph{Cartier} if there are one model $Z$ and one class $a_Z$ such that its class on every model $Z'$ dominating $Z$ is the pullback of $a_Z$.
\end{definition}

When the models are K\"ahler, the standard $\partial\bar\partial$ lemma identifies this cohomology with the real $(1,1)$-cohomology used in \cite[section \emph{The notion of b-divisors}]{SGPT}. The definition then agrees with the source definition.

For the trace data fixed at the beginning of the chapter, the common majorant supplies positive representatives of the residual classes.

\begin{proposition}\label{prop:trace-bott-chern-bclass}\index{Bott--Chern class!residual trace}
If the original ambient background $\theta$ is closed, the classes
\[
 \mathbb A_Z(u;Y,\theta)
 =\{\pi^*(\theta|_Y)-S_Z(u;Y)\}_{\mathrm{BC}}
\]
form a pseudoeffective Bott--Chern b-divisor class. They depend only on the ambient multiplier-ideal class of $u$ and the Bott--Chern class of $\theta|_Y$.
\end{proposition}
\begin{proof}
By \cref{prop:original-background-trace-majorant}, there is one $v$ on $Y$ such that
\[
 R_Z=\pi^*(\theta|_Y)+\ddc(v\circ\pi)-S_Z(u;Y)\geq0
\]
on every model. Since $\theta$ and $S_Z$ are closed, $R_Z$ is closed. Its class is exactly the displayed one, proving pseudoeffectivity. Compatibility follows from \cref{thm:trace-divisorial-model-compatibility} and $g_*g^*(\theta_Z)=\theta_Z$. The trace data is unchanged by ambient multiplier-ideal equivalence, and replacing $\theta|_Y$ by a smooth $\ddc$ perturbation changes no class. These facts prove independence.
\end{proof}

\Needspace{7\baselineskip}
If $\theta$ is not closed, the same expression $R_Z$ is still a positive current, but
\[
 \mathrm{d}R_Z=\pi^*(\mathrm{d}\theta|_Y).
\]
It therefore need not define a class in ordinary Bott--Chern cohomology. The larger classes with a specified non-closed derivative will be treated in \cref{chap:nonclosed-bclasses}.

\Needspace{12\baselineskip}
\section{The b-divisor class of a closed positive current}

For a closed positive current, the same construction uses its own local potentials in place of restricted ambient ideals. The trace coefficients are then its divisorial Lelong numbers, and divisorial removal gives the positive residual currents directly.

Let $T$ be a closed positive $(1,1)$-current on a connected compact complex manifold $Y$, and write $T=\beta+\ddc w$ as above. On a compact modification $\pi:Z\to Y$, local psh potentials define the positive pullback $\pi^*T$: dominance prevents their compositions from being identically $-\infty$. With $E$ ranging over prime divisors on $Z$, set
\[
 N_Z(T)=\sum_E\nu(w\circ\pi,E)[E],\qquad
 \operatorname{Reg}_Z(T)=\pi^*T-N_Z(T).
\]
Divisorial removal makes the second current positive and closed. The following lemma identifies the coefficients with the trace construction.

\begin{lemma}\label{lem:self-trace-is-siu}
For every quasi-plurisubharmonic function $w$ on $Y$ and every prime divisor $E$ on a smooth compact modification $\pi:Z\to Y$, the self-trace in \eqref{eq:trace-profile} satisfies
\[
 \tau_{Y,w}(E)=\nu(w\circ\pi,E).
\]
In particular, for the potential of the closed positive current $T$, one has $N_Z(T)=S_Z(w;Y)$.
\end{lemma}
\begin{proof}
A genuine quasi-plurisubharmonic function has zero Lelong number almost everywhere, so its self-trace satisfies the required generic hypothesis. For a fixed divisor $E$ on a compact modification, apply \cref{lem:bergman-divisorial-error} on a coordinate neighborhood of its center. With $\nu_E=\nu(w\circ\pi,E)$ and $A_E$ the log discrepancy, it gives
\[
 \max\{\nu_E-A_E/m,0\}
 \leq\frac1m\operatorname{ord}_E\mathcal I(mw)
 \leq\nu_E.
\]
The integer limit defining $\tau_{Y,w}(E)$ is therefore $\nu_E$. Equality of all coefficients identifies the two divisorial currents.
\end{proof}

\Needspace{13\baselineskip}
\begin{theorem}\label{thm:closed-current-bclass}
The family
\[
 \mathbb D(T)_Z=\{\operatorname{Reg}_Z(T)\}_{\mathrm{BC}}
\]
is a pseudoeffective Bott--Chern b-divisor class. For closed positive currents $T_1,T_2$ and $a\geq0$,
\[
 \mathbb D(T_1+T_2)=\mathbb D(T_1)+\mathbb D(T_2),
 \qquad \mathbb D(aT_1)=a\mathbb D(T_1).
\]
On a K\"ahler manifold, with K\"ahler models, it agrees with the b-divisor class of the same current in \cite[subsection \emph{The b-divisors of currents}]{SGPT}.
\end{theorem}
\begin{proof}
The preceding lemma identifies $N_Z(T)$ with the already constructed compatible divisorial currents. For $g:Z'\to Z$, we have $g_*N_{Z'}(T)=N_Z(T)$ and $g_*g^*(\pi^*T)=\pi^*T$. For the latter identity, write the current as a smooth form plus $\ddc$ of a potential; degree-one change of variables applies both to the smooth form and to the locally integrable pulled-back potential, and pushforward commutes with $\ddc$. Therefore
\[
 g_*\operatorname{Reg}_{Z'}(T)=\operatorname{Reg}_Z(T).
\]
Taking classes proves compatibility, and positivity proves pseudoeffectivity. Generic Lelong numbers add under addition of positive currents and scale by nonnegative constants. Hence their divisorial parts, and then their regular parts, have the same properties. This proves the two identities. In the K\"ahler case both the regular-current construction and the identified cohomology classes are exactly those in the source definition.
\end{proof}

\section{Nef Cartier realizations}\label{sec:nef-cartier-realizations}

If $T$ has analytic singularities, a resolution removes them in a finite divisor and leaves a current with bounded potentials. This remainder acquires no new divisorial part on higher models, so its class determines the entire family there. Smooth regularization of its bounded potential makes the realizing class nef.

\index{Nef class!definition}A class $a\in H^{1,1}_{\mathrm{BC}}(Z,\mathbb R)$ is \emph{nef} if, for one Hermitian metric $\gamma$, it has a smooth representative $\beta_\epsilon\geq-\epsilon\gamma$ for every $\epsilon>0$. Comparison of two fixed Hermitian metrics shows that this definition is independent of the metric. A Cartier b-divisor class is called nef here if it has a realization by a nef class.

\Needspace{13\baselineskip}
\begin{lemma}\label{lem:bounded-closed-potential-nef}
On a connected compact complex manifold $Z$, let $\beta$ be a smooth real $(1,1)$-form and $v$ a bounded quasi-plurisubharmonic function. If $R=\beta+\ddc v\geq0$ is closed, then $\{R\}_{\mathrm{BC}}$ is nef.
\end{lemma}
\begin{proof}
Since $R$ is closed, $\beta$ is closed. Fix a Hermitian metric $\gamma$. By \cref{prop:bounded-almost-regularization}, there are smooth $v_j$ and $\epsilon_j\downarrow0$ such that
\[
 \beta+\ddc v_j\geq-\epsilon_j\gamma.
\]
These are closed representatives of $\{R\}_{\mathrm{BC}}$. Given any $\epsilon>0$, choose $j$ with $\epsilon_j\leq\epsilon$. This is the defining nef estimate.
\end{proof}

\Needspace{8\baselineskip}
\begin{theorem}\label{thm:analytic-current-nef-cartier}\index{Nef class!Cartier realization}
If the closed positive current $T$ has analytic singularities, then $\mathbb D(T)$ is Cartier and has a nef realization on a smooth compact modification.
\end{theorem}
\begin{proof}
Resolve the analytic singularity ideal, using the log-resolution theorem in \cite[section on analytic singularities, theorem on existence of log resolutions]{SGPT}. On one smooth compact model $\pi:Z\to Y$,
\[
 \pi^*T=[D]+R,
\]
where $D$ is a finite effective rational divisor and $R$ is positive, closed, and has locally bounded potentials. To obtain a global bounded potential, choose smooth metrics on the divisor bundles. Subtract the corresponding logarithmic squared norms from $w\circ\pi$. The remainder $v$ is globally bounded, and
\[
 R=\beta_Z+\ddc v
\]
for a smooth closed $\beta_Z$. The lemma makes its class nef.

On every further compact modification $g:Z'\to Z$, the current $g^*R$ still has bounded local potentials and consequently zero Lelong numbers. The entire divisorial part of $g^*\pi^*T$ is therefore $[g^*D]$, and its regular part is $g^*R$. It follows that $\mathbb D(T)_{Z'}=g^*\{R\}_{\mathrm{BC}}$. Thus the pair $(Z,\{R\}_{\mathrm{BC}})$ is the required nef Cartier realization.
\end{proof}

For a general current, analytic regularization gives closed approximants bounded below by $-\epsilon_m\gamma$. Adding $\epsilon_m\gamma$ makes them positive but changes their exterior derivative when $\gamma$ is not closed. The approximation then belongs to the larger spaces of \cref{chap:nonclosed-bclasses}; it does not by itself approximate the original Bott--Chern b-divisor class by nef Cartier classes.
